\chapter{Correlation Analyses and Normalization in R}
\label{chap:correlation-normalization}

\chapterguide
  {You should be comfortable loading datasets, selecting columns, installing packages, and creating visualizations.}
  {compute Pearson correlation using \texttt{cor()} and \texttt{cor.test()}; visualize correlation matrices using the packages shown in the handout; construct correlation heatmaps; and compare log transformation, standard scaling, and min--max scaling using the supplied examples.}

\section{Pearson correlation with \texttt{cor()}}

The handout begins with two vectors:

\begin{lstlisting}[style=dsr]
x <- c(1, 2, 3, 4, 5, 6, 7)
y <- c(1, 3, 6, 2, 7, 4, 5)

result <- cor(x, y, method = "pearson")
cat("Pearson correlation coefficient is:", result)
\end{lstlisting}

The purpose is to obtain a single Pearson correlation coefficient describing the linear association between the paired values.

\section{Pearson correlation with \texttt{cor.test()}}

The next example uses the same vectors but requests a full test object:

\begin{lstlisting}[style=dsr]
result <- cor.test(x, y, method = "pearson")
print(result)
\end{lstlisting}

The handout therefore distinguishes a direct coefficient calculation from a more detailed statistical-test output.

\section{Correlation matrix with \texttt{corrplot}}

Using \texttt{mtcars}, the source computes a correlation matrix and visualizes it:

\begin{lstlisting}[style=dsr]
install.packages("corrplot")
install.packages("RColorBrewer")

library(corrplot)
library(RColorBrewer)

M <- cor(mtcars)
?corrplot

corrplot(M, type = "upper")
corrplot(M, type = "upper", order = "hclust")
corrplot(
  M,
  type = "upper",
  order = "hclust",
  col = brewer.pal(n = 8, name = "RdYlBu")
)
\end{lstlisting}

The progression adds hierarchical ordering and then a color palette.

\section{Building a heatmap-ready correlation table}

Lab 9a next uses the built-in \texttt{environmental} data from the \texttt{lattice} package:

\begin{lstlisting}[style=dsr]
install.packages("lattice")
library(lattice)

data <- environmental
head(data)

corr_mat <- round(cor(data), 2)
head(corr_mat)
\end{lstlisting}

The matrix is reordered using a distance derived from correlation values and hierarchical clustering:

\begin{lstlisting}[style=dsr]
dist <- as.dist((1 - corr_mat) / 2)
hc <- hclust(dist)
corr_mat <- corr_mat[hc$order, hc$order]
\end{lstlisting}

Then \texttt{reshape2::melt()} converts the matrix into a longer form:

\begin{lstlisting}[style=dsr]
install.packages("reshape2")
library(reshape2)

melted_corr_mat <- melt(corr_mat)
head(melted_corr_mat)
\end{lstlisting}

This transformed table is the input for the following \texttt{ggplot2} heatmap.

\section{Correlation heatmaps}

\subsection{Using \texttt{ggplot2}}

\begin{lstlisting}[style=dsr]
install.packages("ggplot2")
library(ggplot2)

ggplot(
  data = melted_corr_mat,
  aes(x = Var1, y = Var2, fill = value)
) +
  geom_tile() +
  geom_text(
    aes(Var2, Var1, label = value),
    color = "white",
    size = 4
  )
\end{lstlisting}

The heatmap encodes each correlation value as a tile and overlays the numeric value as text.

\subsection{Using \texttt{heatmaply}}

\begin{lstlisting}[style=dsr]
install.packages("heatmaply")
library(heatmaply)

heatmaply_cor(
  x = cor(data),
  xlab = "Features",
  ylab = "Features",
  k_col = 2,
  k_row = 2
)
\end{lstlisting}

\subsection{Using \texttt{ggcorrplot}}

The handout's package-installation line and library name do not match exactly:

\begin{lstlisting}[style=dsr]
# Printed package name in the handout:
install.packages("ggcorplot")

# Printed library/function name:
library(ggcorrplot)
ggcorrplot::ggcorrplot(cor(data))
\end{lstlisting}

\begin{pitfall}
Lab 9a prints \texttt{install.packages("ggcorplot")} but immediately uses \texttt{library(ggcorrplot)} and \texttt{ggcorrplot::ggcorrplot()}. These notes preserve the mismatch as a source caution. Verify the package name in your R environment before installation rather than assuming the two spellings are equivalent.
\end{pitfall}

\section{Normalization by log transformation}

The supplied numeric vector is:

\begin{lstlisting}[style=dsr]
mydata <- c(244, 753, 596, 645, 874, 141, 639, 465, 999, 654)

scaled_data <- log(mydata)
print(scaled_data)
\end{lstlisting}

The lab labels this as normalization with log transformation and asks students to inspect how the values change.

\section{Standard scaling}

The handout uses \texttt{scale()} and converts the result to a data frame:

\begin{lstlisting}[style=dsr]
scaled_data2 <- as.data.frame(scale(mydata))
print(scaled_data2)
\end{lstlisting}

This gives a second transformed representation of exactly the same source values.

\section{Min--max scaling}

The third method uses \texttt{caret::preProcess()} with the \texttt{"range"} method:

\begin{lstlisting}[style=dsr]
install.packages("caret")
library(caret)

minmax <- preProcess(
  as.data.frame(mydata),
  method = c("range")
)

scaled_data3 <- predict(
  minmax,
  as.data.frame(mydata)
)
print(scaled_data3)
\end{lstlisting}

\section{Comparing summaries}

The final Lab 9a step is to compare summary statistics for the raw and transformed data:

\begin{lstlisting}[style=dsr]
summary(mydata)
summary(scaled_data1)
summary(scaled_data2)
summary(scaled_data3)
\end{lstlisting}

\begin{pitfall}
The log-transformation code stores its result in \texttt{scaled\_data}, but the final comparison calls \texttt{summary(scaled\_data1)}. The source does not define \texttt{scaled\_data1}. To execute the intended comparison, either use \texttt{summary(scaled\_data)} or explicitly assign \texttt{scaled\_data1 <- scaled\_data} before the summary calls. This is an explicit implementation note, not a silent source correction.
\end{pitfall}

\begin{dsfigure}
\begin{tikzpicture}[
  box/.style={draw=DSBlue,fill=DSPaleBlue,rounded corners=2mm,text width=3.35cm,minimum height=1.1cm,align=center,font=\small\bfseries\color{DSNavy}},
  arr/.style={-{Stealth[length=2mm]},thick,draw=DSTeal}
]
\node[box] (raw) {Raw numeric data};
\node[box,right=1.0cm of raw] (log) {Log transformation};
\node[box,below=0.8cm of log] (std) {Standard scaling};
\node[box,above=0.8cm of log] (mm) {Min--max scaling};
\draw[arr] (raw) -- (log);
\draw[arr] (raw.east) -- ++(0.55,0) |- (std.west);
\draw[arr] (raw.east) -- ++(0.55,0) |- (mm.west);
\end{tikzpicture}
\end{dsfigure}

The three branches represent the three transformation methods explicitly requested in Lab 9a; the final exercise compares their summary output against the raw data.

\begin{quickcheck}
\begin{enumerate}
  \item What is the difference in output intent between \texttt{cor()} and \texttt{cor.test()} in Lab 9a?
  \item Which dataset is used for the first \texttt{corrplot()} example?
  \item Why is \texttt{melt()} applied to the reordered correlation matrix before the \texttt{ggplot2} heatmap?
  \item Which function performs the standard scaling in the handout?
  \item Which object-name mismatch appears in the final summary comparison?
\end{enumerate}
\end{quickcheck}

\begin{chapterrecap}
\begin{itemize}
  \item Pearson correlation is computed directly with \texttt{cor()} and examined more fully with \texttt{cor.test()}.
  \item The source demonstrates several correlation visualizations: \texttt{corrplot}, a \texttt{ggplot2} tile heatmap, \texttt{heatmaply}, and \texttt{ggcorrplot}.
  \item Correlation matrices can be reordered before plotting to group similar relationships visually.
  \item The normalization section compares log transformation, standard scaling, and min--max scaling on the same data.
  \item Lab 9a contains a few naming/package inconsistencies that should be checked explicitly when reproducing the code.
\end{itemize}
\end{chapterrecap}
